% GENERATED FILE. Edit canonical lessons, then run npm run generate:books.
% canonical-source-hash: fnv1a64:9b4ddebbd5f41f48
% canonical-lessons: MR-C556-anukul, MR-C556-pratikul, MR-C556-ghol, MR-C556-bikat, MR-C556-sambhram

\chapter{Favourable or Not}
\label{ch:mr-favourable-or-not}
\hlchaptermodality{556}

\begin{chapteropening}
\textbf{By the end of this chapter:} \emph{I can say whether conditions are favourable and talk about muddles.}

Complete the last lesson of chapter 556: I can say whether conditions are favourable and talk about muddles.
\end{chapteropening}

\section[\texorpdfstring{\mr{अनुकूल} (anukūl)}{अनुकूल (anukūl)}]{\mr{अनुकूल} (anukūl) — favourable}

\label{lesson:MR-C556-anukul}



\begin{quote}
Before the new one: say the Marathi for a rush, then the Marathi for a scramble.
\end{quote}

\subsection*{You'll want to know: \mr{अनुकूल}}
\textbf{\mr{अनुकूल}} — \emph{anukūl} — \textquotedblleft{}favourable\textquotedblright{}.

\begin{quote}
\mr{हवामान} \mr{अनुकूल} \mr{असेल} \mr{तर} \mr{आम्ही} \mr{उद्या} \mr{निघू}.
\end{quote}

\begin{quote}
\emph{havāmān anukūl asel tar āmhī udyā nighū.}
\end{quote}

\begin{quote}
If the weather is favourable we will set out tomorrow.
\end{quote}

\begin{cousinweb}
\textbf{\mr{अनुकूल}} (\emph{anukūl}) is a tatsama, taken from Sanskrit \emph{anukūla}, \textquotedblleft{}following the bank or slope\textquotedblright{} (Monier-Williams), so going with the stream.
\end{cousinweb}

\begin{grammarlens}[title={What you've built}]
One more word for saying what happens if.
\end{grammarlens}

\subsection*{Your turn}
\begin{itemize}
\raggedright
  \item \emph{Say it:} \emph{anukūl}
  \item \emph{Say it:} \emph{anukūl}, once more
  \item \emph{Say it:} \emph{havāmān anukūl asel tar āmhī udyā nighū.} — If the weather is favourable we will set out tomorrow.
  \item \emph{Recall:} say the Marathi for a rush, then the Marathi for a scramble, then say \emph{anukūl} again
\end{itemize}

\begin{tcolorbox}[breakable,colback=teal!4,colframe=teal!35!black,title={Before you move on}]
What does \mr{अनुकूल} mean? (\textquotedblleft{}Favourable\textquotedblright{}.) Say it once more.
\end{tcolorbox}
\section[\texorpdfstring{\mr{प्रतिकूल} (pratikūl)}{प्रतिकूल (pratikūl)}]{\mr{प्रतिकूल} (pratikūl) — unfavourable; adverse}

\label{lesson:MR-C556-pratikul}



\begin{quote}
Before the new one: say the Marathi for a scramble, then the Marathi for favourable.
\end{quote}

\subsection*{You'll want to know: \mr{प्रतिकूल}}
\textbf{\mr{प्रतिकूल}} — \emph{pratikūl} — \textquotedblleft{}unfavourable; adverse\textquotedblright{}.

\begin{quote}
\mr{प्रतिकूल} \mr{परिस्थितीतही} \mr{ती} \mr{शिकत} \mr{राहिली}.
\end{quote}

\begin{quote}
\emph{pratikūl paristhitīthī tī śikat rāhilī.}
\end{quote}

\begin{quote}
Even in adverse circumstances she kept on studying.
\end{quote}

\begin{cousinweb}
\textbf{\mr{प्रतिकूल}} (\emph{pratikūl}) is a tatsama, taken from Sanskrit \emph{pratikūla}, \textquotedblleft{}against the bank, adverse\textquotedblright{} (Monier-Williams).
\end{cousinweb}

\begin{grammarlens}[title={What you've built}]
One more word for saying what happens if.
\end{grammarlens}

\subsection*{Your turn}
\begin{itemize}
\raggedright
  \item \emph{Say it:} \emph{pratikūl}
  \item \emph{Say it:} \emph{pratikūl}, once more
  \item \emph{Say it:} \emph{pratikūl paristhitīthī tī śikat rāhilī.} — Even in adverse circumstances she kept on studying.
  \item \emph{Recall:} say the Marathi for a scramble, then the Marathi for favourable, then say \emph{pratikūl} again
\end{itemize}

\begin{tcolorbox}[breakable,colback=teal!4,colframe=teal!35!black,title={Before you move on}]
What does \mr{प्रतिकूल} mean? (\textquotedblleft{}Unfavourable; adverse\textquotedblright{}.) Say it once more.
\end{tcolorbox}
\section[\texorpdfstring{\mr{घोळ} (ghoḷ)}{घोळ (ghoḷ)}]{\mr{घोळ} (ghoḷ) — a muddle; a mess}

\label{lesson:MR-C556-ghol}



\begin{quote}
Before the new one: say the Marathi for favourable, then the Marathi for unfavourable.
\end{quote}

\subsection*{You'll want to know: \mr{घोळ}}
\textbf{\mr{घोळ}} — \emph{ghoḷ} — \textquotedblleft{}a muddle; a mess\textquotedblright{}.

\begin{quote}
\mr{तिकिटांच्या} \mr{आरक्षणात} \mr{मोठा} \mr{घोळ} \mr{झाला}.
\end{quote}

\begin{quote}
\emph{tikiṭāñcyā ārakṣaṇāt moṭhā ghoḷ jhālā.}
\end{quote}

\begin{quote}
There was a big muddle in the ticket booking.
\end{quote}

Masculine.

\mr{घोळ} \mr{होणे}, things get into a muddle.

\begin{cousinweb}
\textbf{\mr{घोळ}} (\emph{ghoḷ}) is built on \mr{घोळणे}, \textquotedblleft{}to stir about\textquotedblright{}; CDIAL lists Marathi \emph{ghoḷṇẽ} under Sanskrit \emph{ghōlayati}, \textquotedblleft{}mixes, stirs together\textquotedblright{} (CDIAL 4526).
\end{cousinweb}

\begin{grammarlens}[title={What you've built}]
One more word for saying what happens if.
\end{grammarlens}

\subsection*{Your turn}
\begin{itemize}
\raggedright
  \item \emph{Say it:} \emph{ghoḷ}
  \item \emph{Say it:} \emph{ghoḷ}, once more
  \item \emph{Say it:} \emph{tikiṭāñcyā ārakṣaṇāt moṭhā ghoḷ jhālā.} — There was a big muddle in the ticket booking.
  \item \emph{Recall:} say the Marathi for favourable, then the Marathi for unfavourable, then say \emph{ghoḷ} again
\end{itemize}

\begin{tcolorbox}[breakable,colback=teal!4,colframe=teal!35!black,title={Before you move on}]
What does \mr{घोळ} mean? (\textquotedblleft{}A muddle; a mess\textquotedblright{}.) Say it once more.
\end{tcolorbox}
\section[\texorpdfstring{\mr{बिकट} (bikaṭ)}{बिकट (bikaṭ)}]{\mr{बिकट} (bikaṭ) — dire; critical}

\label{lesson:MR-C556-bikat}



\begin{quote}
Before the new one: say the Marathi for unfavourable, then the Marathi for a muddle.
\end{quote}

\subsection*{You'll want to know: \mr{बिकट}}
\textbf{\mr{बिकट}} — \emph{bikaṭ} — \textquotedblleft{}dire; critical\textquotedblright{}.

\begin{quote}
\mr{पुरामुळे} \mr{गावातली} \mr{परिस्थिती} \mr{बिकट} \mr{झाली}.
\end{quote}

\begin{quote}
\emph{purāmuḷe gāvātlī paristhitī bikaṭ jhālī.}
\end{quote}

\begin{quote}
Because of the flood, things in the village became dire.
\end{quote}

Mostly of a situation or a road: \mr{बिकट} \mr{परिस्थिती}, \mr{बिकट} \mr{वाट}.

\begin{cousinweb}
\textbf{\mr{बिकट}} (\emph{bikaṭ}) is usually traced to Sanskrit \emph{vikaṭa}, \textquotedblleft{}dreadful, monstrous, of unusual size\textquotedblright{} (Monier-Williams), with the v- said as b-; CDIAL lists no Marathi \mr{बिकट}, so the path is not settled.
\end{cousinweb}

\begin{grammarlens}[title={What you've built}]
One more word for saying what happens if.
\end{grammarlens}

\subsection*{Your turn}
\begin{itemize}
\raggedright
  \item \emph{Say it:} \emph{bikaṭ}
  \item \emph{Say it:} \emph{bikaṭ}, once more
  \item \emph{Say it:} \emph{purāmuḷe gāvātlī paristhitī bikaṭ jhālī.} — Because of the flood, things in the village became dire.
  \item \emph{Recall:} say the Marathi for unfavourable, then the Marathi for a muddle, then say \emph{bikaṭ} again
\end{itemize}

\begin{tcolorbox}[breakable,colback=teal!4,colframe=teal!35!black,title={Before you move on}]
What does \mr{बिकट} mean? (\textquotedblleft{}Dire; critical\textquotedblright{}.) Say it once more.
\end{tcolorbox}
\section[\texorpdfstring{\mr{संभ्रम} (sambhram)}{संभ्रम (sambhram)}]{\mr{संभ्रम} (sambhram) — confusion; doubt}

\label{lesson:MR-C556-sambhram}



\begin{quote}
Before the new one: say the Marathi for a muddle, then the Marathi for dire.
\end{quote}

\subsection*{You'll want to know: \mr{संभ्रम}}
\textbf{\mr{संभ्रम}} — \emph{sambhram} — \textquotedblleft{}confusion; doubt\textquotedblright{}.

\begin{quote}
\mr{दोन} \mr{वेगळ्या} \mr{सूचनांमुळे} \mr{सगळ्यांच्या} \mr{मनात} \mr{संभ्रम} \mr{होता}.
\end{quote}

\begin{quote}
\emph{don vegaḷyā sūcanāmmuḷe sagḷyāñcyā manāt sambhram hotā.}
\end{quote}

\begin{quote}
Because of two different instructions, everyone was confused.
\end{quote}

Masculine.

\begin{cousinweb}
\textbf{\mr{संभ्रम}} (\emph{sambhram}) is a tatsama, taken from Sanskrit \emph{sambhrama}, \textquotedblleft{}whirling round, confusion\textquotedblright{} (Monier-Williams).
\end{cousinweb}

\begin{grammarlens}[title={What you've built}]
One more word for saying what happens if.
\end{grammarlens}

\subsection*{Your turn}
\begin{itemize}
\raggedright
  \item \emph{Say it:} \emph{sambhram}
  \item \emph{Say it:} \emph{sambhram}, once more
  \item \emph{Say it:} \emph{don vegaḷyā sūcanāmmuḷe sagḷyāñcyā manāt sambhram hotā.} — Because of two different instructions, everyone was confused.
  \item \emph{Recall:} say the Marathi for a muddle, then the Marathi for dire, then say \emph{sambhram} again
\end{itemize}

\begin{tcolorbox}[breakable,colback=teal!4,colframe=teal!35!black,title={Before you move on}]
What does \mr{संभ्रम} mean? (\textquotedblleft{}Confusion; doubt\textquotedblright{}.) Say it once more.
\end{tcolorbox}
